\section{An improved explicit polynomial exponent}
\label{sec:two-scale-optimized-constants}

This section changes only the parameters of the two-scale proof.
The structural argument is unchanged.
It also uses paired consecutive spectral powers to shorten the counted-cycle interval.
The original audited parameter choice remains available separately.

\begin{theorem}[An explicit exponent]
\label{thm:two-scale-explicit-exponent}
Put
\[
 \kappa_* = 2^{-1050}.
\]
For all sufficiently large ranks $r$,
\[
 \psi(r)\le2^{r-1}+(1+o(1))2^r r^{-\kappa_*}.
\]
More generally, the polynomial forest bound with this exponent holds whenever
\[
 |P|\le\frac{r}{112000\log_2r},\qquad
 \operatorname{comp}(\Gamma[P])\le\frac{\log_2r}{100},
\]
under the degree hypotheses of Theorem~\ref{thm:two-scale-critical-marking}.
\end{theorem}

\begin{proof}
Set
\[
 x=\log_2r,\qquad A=112,\qquad t=1024,\qquad L=\lfloor112x\rfloor,
\]
\[
 D_0=\lfloor12x/25\rfloor,\qquad
 D_1=\lfloor49x/50\rfloor,\qquad
 s_j=\lfloor(D_j-1)/2\rfloor.
\]
Assume that the forest probability exceeds $1/2+r^{-\kappa_*}$.
Choose the nested families $\mathcal Q_0\subseteq\mathcal Q_1$ as before.
Their vertex unions satisfy
\[
 |S_0|=O(x^3r^{12/25}),\qquad
 |S_1|=O(x^3r^{49/50}).
\]
Both are $o(r/x)$.
The graphs $J_1\subseteq J_0$ have the required girths, and their orders are $2r+o(r)$.

We first check the cycle supply with the new value of $A$.
Let $a$ be the smallest even integer at least $2x+10\log_2x$.
Let $b$ be the largest odd integer at most $L$.
The non-backtracking eigenvalues have modulus at most two.
For a negative real eigenvalue $\lambda=-2u$, with $0\le u\le1$, an even power $j$ and its successor contribute
\[
 \frac{u^j}{2j}-\frac{u^{j+1}}{2(j+1)}\ge0
\]
to the normalized trace sum.
Thus pairing every even power with the following odd power removes any loss from negative real eigenvalues.

The marked-set hypothesis gives
\[
 \frac{Ld_1}{2v_1}\le\frac3{4000}+o(1)<\frac1{500}.
\]
The positive eigenvalue at least $2-d_1/v_1$ consequently contributes at least
\[
 \left(1-\frac1{500}\right)\sum_{j=a}^b\frac1{2j}.
\]
The nonreal contribution has magnitude at most
\[
 3v_1\sum_{j\ge a}\frac{2^{-j/2}}{2j}=O(x^{-6})=o(1),
\]
since $2^{-a/2}\le r^{-1}x^{-5}$.
The repeated-walk factor is $1+O(L^2 2^{-s_1})=1+o(1)$.
Therefore
\begin{equation}\label{eq:optimized-initial-mass}
 \mu_1\ge\left(1-\frac1{500}\right)\frac{\ln56}{2}-o(1)
 >\frac{401}{200}=2.005
\end{equation}
for sufficiently large $r$.
The strict numerical inequality has a margin greater than $0.003$.

Next, check the uniform link bound at the smaller girth scale.
The marked component term is at most $x/100$.
Also
\[
 \frac{2L}{t}\le\frac7{32}x,
 \qquad
 8t2^t\ln(2r^{\kappa_*})
 =\frac{\ln2}{8192}x+O(1).
\]
The remaining tree-link term is bounded independently of $r$.
Since
\[
 \frac1{100}+\frac7{32}+\frac{\ln2}{8192}<\frac{47}{200}=0.235,
\]
all link counts are at most $47x/200$ for large $r$.
The tree argument's condition $2t+D_0/2\le D_0$ also holds eventually.
Because $s_0=6x/25+O(1)$, the exceptional contribution is at most
\[
 a_C\,2^{47x/200}L^3 2^{-s_0}
 =O(a_C x^3r^{-1/200})=o(a_C).
\]
Thus both inverse moment bounds remain valid.

Use one connected component of $\Gamma[P]$ as the anchor for far-pair antisymmetry.
The weighted main-component partner mass for some cycle is at least
$\mu_1/2-1-o(1)$.
By \eqref{eq:optimized-initial-mass}, this is at least $1/1000$ for large $r$.
The same maximal-disjoint-family argument supplies a nested chain with
\[
 T\ge\frac{2^{s_1}}{1000L^3}
 =\Omega(r^{49/100}/x^3).
\]

Group the chain into blocks using $q=\lceil20000Ax\rceil+2$ as before.
There are $\Omega(r^{49/100}/x^4)$ blocks.
The blocks have boundary at most $2L$ and volume at least $2000Ax^2$.
Discard blocks meeting a marked component, a neighbor of the anchor cycle, or $S_0$ and its neighbors.
The total loss is at most
\[
 \operatorname{comp}(\Gamma[P])+L+4|S_0|
 =O(x^3r^{12/25})
 =o(r^{49/100}/x^4).
\]
Every remaining block has zero deficit in $J_0$.

The multiple-mode spectral test, still using $\theta=1/(12x)$ and even lengths between $6x$ and $24x$, now gives
\[
 \mu_0=\Omega(r^{49/100}/x^4).
\]
Those lengths lie below $L=\lfloor112x\rfloor$.
The inverse moment estimate gives
\[
 \mu_0\le\frac{1+o(1)}4\,r^{\kappa_*},
\]
a contradiction since $\kappa_*<12/25$.
This proves the forest estimate.
The direct cubic-extremizer transfer and the same geometric recurrence prove the cycle estimate.
\end{proof}

The exponent remains too small for a useful finite-rank estimate.
This calculation gives a substantially simpler explicit constant without changing the theorem's scope or proof structure.
